% Options for packages loaded elsewhere
\PassOptionsToPackage{unicode}{hyperref}
\PassOptionsToPackage{hyphens}{url}
\documentclass[
  ignorenonframetext,
  aspectratio=169,
]{beamer}
\newif\ifbibliography
\usepackage{pgfpages}
\setbeamertemplate{caption}[numbered]
\setbeamertemplate{caption label separator}{: }
\setbeamercolor{caption name}{fg=normal text.fg}
\beamertemplatenavigationsymbolsempty
% remove section numbering
\setbeamertemplate{part page}{
  \centering
  \begin{beamercolorbox}[sep=16pt,center]{part title}
    \usebeamerfont{part title}\insertpart\par
  \end{beamercolorbox}
}
\setbeamertemplate{section page}{
  \centering
  \begin{beamercolorbox}[sep=12pt,center]{section title}
    \usebeamerfont{section title}\insertsection\par
  \end{beamercolorbox}
}
\setbeamertemplate{subsection page}{
  \centering
  \begin{beamercolorbox}[sep=8pt,center]{subsection title}
    \usebeamerfont{subsection title}\insertsubsection\par
  \end{beamercolorbox}
}
% Prevent slide breaks in the middle of a paragraph
\widowpenalties 1 10000
\raggedbottom
\AtBeginPart{
  \frame{\partpage}
}
\AtBeginSection{
  \ifbibliography
  \else
    \frame{\sectionpage}
  \fi
}
\AtBeginSubsection{
  \frame{\subsectionpage}
}
\usepackage{iftex}
\ifPDFTeX
  \usepackage[T1]{fontenc}
  \usepackage[utf8]{inputenc}
  \usepackage{textcomp} % provide euro and other symbols
\else % if luatex or xetex
  \usepackage{unicode-math} % this also loads fontspec
  \defaultfontfeatures{Scale=MatchLowercase}
  \defaultfontfeatures[\rmfamily]{Ligatures=TeX,Scale=1}
\fi
\usepackage{lmodern}
\ifPDFTeX\else
  % xetex/luatex font selection
\fi
% Use upquote if available, for straight quotes in verbatim environments
\IfFileExists{upquote.sty}{\usepackage{upquote}}{}
\IfFileExists{microtype.sty}{% use microtype if available
  \usepackage[]{microtype}
  \UseMicrotypeSet[protrusion]{basicmath} % disable protrusion for tt fonts
}{}
\makeatletter
\@ifundefined{KOMAClassName}{% if non-KOMA class
  \IfFileExists{parskip.sty}{%
    \usepackage{parskip}
  }{% else
    \setlength{\parindent}{0pt}
    \setlength{\parskip}{6pt plus 2pt minus 1pt}}
}{% if KOMA class
  \KOMAoptions{parskip=half}}
\makeatother
\usepackage{longtable,booktabs,array}
\usepackage{caption}
\captionsetup[table]{skip=6pt}
\newcounter{none} % for unnumbered tables
\usepackage{calc} % for calculating minipage widths
\usepackage{caption}
% Make caption package work with longtable
\makeatletter
\def\fnum@table{\tablename~\thetable}
\makeatother
\setlength{\emergencystretch}{3em} % prevent overfull lines
\providecommand{\tightlist}{%
  \setlength{\itemsep}{0pt}\setlength{\parskip}{0pt}}
% Auto-generated by infrastructure/rendering/_pdf_latex_helpers.py::extract_math_font_preamble.
% Minimal header passed to Pandoc via -H so Beamer slides pick up the same math font
% as the combined-PDF path. Adjust the manuscript's preamble.md to change the font globally.
\usepackage{unicode-math}
\setmathfont{latinmodern-math.otf}

% Manuscript macro/package fallbacks (newcommand -> providecommand).
% Core mathematics
\usepackage{amsmath}
\usepackage{amssymb}
% Document layout
% Tables
\usepackage{booktabs}
% Caption styling (booktabs already loaded above). Small caption font with a
% bold label ("Figure 1", "Table 2") gives the math/figure-heavy layout a
% consistent, journal-like caption rhythm without fighting pandoc-crossref —
% pandoc emits standard \caption{}, and caption/captionsetup only restyle it.
% ── Theorem environments ─────────────────────────────────────────────────
% amsthm is loaded above. The FedGVI<->Friston bridge is stated as a small
% number of theorems/definitions/lemmas/propositions/corollaries; sharing one
% counter (definition/lemma/proposition/corollary all step the `theorem`
% counter) gives a single monotone numbering 1, 2, 3, ... across all kinds, so
% authors can reference them as Theorem 1, Definition 2, Corollary 3 without a
% per-kind counter drifting out of order. Plain style for theorem-like results,
% definition style (upright body) for definitions.
% (algorithm2e intentionally not loaded: this manuscript states results as
% numbered amsthm Theorems/Definitions, not algorithm2e pseudocode.)
% Code listings
\usepackage{listings}
% Typography and formatting
% (siunitx intentionally not loaded: no \SI/\num units appear in this manuscript.)
% Cross-references and citations
% (cleveref and titlesec intentionally not loaded: unavailable in this TeX tree
% and not required. Cross-references resolve through pandoc-crossref's own
% \ref-style names — no \cref appears — and section heading styling falls back
% to the pandoc/LaTeX defaults.)
% ── TeX Live Basic-compatible mono font for code listings ────────────
% Latin Modern Mono is bundled with TeX Live Basic and keeps the manuscript
% renderable on minimal TeX installations. Mathematical Unicode coverage is
% provided separately by Latin Modern Math below, so the code font does not
% need a private JuliaMono installation.
%
% Use the TeX font filename rather than the system-family name: the minimal
% TeX Live fontconfig database may not register the OpenType family even though
% kpsewhich can resolve the bundled files.
% Math font for unicode-math: Latin Modern Math (TeX Live) has full BMP
% coverage including U+2223 (\mid), U+226A/226B (\ll/\gg), and the Greek/
% blackboard letters used in equations. Without an explicit \setmathfont,
% unicode-math falls back to lmroman text font which lacks several glyphs
% and emits "Missing character" warnings on every \mid in math mode.

% Slide layout defaults for warning-clean scientific decks.
\usepackage{etoolbox}
\IfFileExists{xurl.sty}{\usepackage{xurl}}{}
\IfFileExists{seqsplit.sty}{\usepackage{seqsplit}}{\newcommand{\seqsplit}[1]{#1}}
\protected\def\breakseq#1{\seqsplit{#1}}
\protected\def\breaktt#1{\begingroup\ttfamily\seqsplit{#1}\endgroup}
\setlength{\emergencystretch}{6em}
\tolerance=5000
\hbadness=10000
\hfuzz=1pt
\setlength{\tabcolsep}{2pt}
\AtBeginEnvironment{longtable}{\tiny\renewcommand{\arraystretch}{0.86}\setlength{\tabcolsep}{1pt}}
\AtBeginEnvironment{tabular}{\tiny\renewcommand{\arraystretch}{0.86}\setlength{\tabcolsep}{1pt}}
\AtBeginEnvironment{equation}{\tiny}
\AtBeginEnvironment{equation*}{\tiny}
\AtBeginEnvironment{align}{\tiny}
\AtBeginEnvironment{align*}{\tiny}
\AtBeginEnvironment{itemize}{\footnotesize}
\AtBeginEnvironment{enumerate}{\footnotesize}
\AtBeginEnvironment{description}{\footnotesize}
\setbeamerfont{caption}{size=\tiny}
\setbeamerfont{caption name}{size=\tiny}
\setbeamerfont{normal text}{size=\small}
\setbeamerfont{frametitle}{size=\small}
\setbeamerfont{section title}{size=\footnotesize}
\setbeamerfont{subsection title}{size=\footnotesize}
\setbeamertemplate{section page}{%
  \centering
  \begin{beamercolorbox}[sep=12pt,center,wd=\paperwidth]{section title}
    \parbox{0.86\paperwidth}{\centering\usebeamerfont{section title}\insertsection\par}
  \end{beamercolorbox}
}
\setbeamertemplate{subsection page}{%
  \centering
  \begin{beamercolorbox}[sep=8pt,center,wd=\paperwidth]{subsection title}
    \parbox{0.86\paperwidth}{\centering\usebeamerfont{subsection title}\insertsubsection\par}
  \end{beamercolorbox}
}
\setlength{\abovecaptionskip}{2pt}
\setlength{\belowcaptionskip}{0pt}

% Natbib and cross-reference fallbacks — slides don't load natbib
% or cleveref, but combined-PDF manuscript prose may emit these
% commands. The fallback renders citations as a bracketed key list
% and cross-references as detokenized labels so slides don't fail on
% undefined control sequences or raw underscores. \providecommand is
% a no-op if packages load later.
\providecommand{\citep}[1]{[#1]}
\providecommand{\citet}[1]{#1}
\providecommand{\citealp}[1]{#1}
\providecommand{\citeauthor}[1]{#1}
\providecommand{\citeyear}[1]{#1}
\providecommand{\cref}[1]{\texttt{\detokenize{#1}}}
\providecommand{\Cref}[1]{\texttt{\detokenize{#1}}}
\providecommand{\crefrange}[2]{\texttt{\detokenize{#1}}--\texttt{\detokenize{#2}}}
\providecommand{\Crefrange}[2]{\texttt{\detokenize{#1}}--\texttt{\detokenize{#2}}}
\providecommand{\paragraph}[1]{\textbf{#1}\ }

% Opt-in accessible presentation profile.
\setbeamerfont{normal text}{size*={20pt}{24pt}}
\setbeamerfont{frametitle}{size*={28pt}{32pt}}
\setbeamerfont{section title}{size*={28pt}{32pt}}
\setbeamerfont{subsection title}{size*={28pt}{32pt}}
\setbeamerfont{caption}{size*={16pt}{19pt}}
\setbeamerfont{caption name}{size*={16pt}{19pt}}
\setbeamerfont{itemize/enumerate body}{size*={20pt}{24pt}}
\setbeamerfont{itemize/enumerate subbody}{size*={20pt}{24pt}}
\setbeamerfont{itemize/enumerate subsubbody}{size*={20pt}{24pt}}
\setbeamerfont{itemize item}{size*={20pt}{24pt}}
\setbeamerfont{itemize subitem}{size*={20pt}{24pt}}
\setbeamerfont{itemize subsubitem}{size*={20pt}{24pt}}
\setbeamerfont{description body}{size*={20pt}{24pt}}
\setbeamerfont{description item}{size*={20pt}{24pt}}
\setbeamerfont{quote}{size*={20pt}{24pt}}
\setbeamerfont{quotation}{size*={20pt}{24pt}}
\AtBeginEnvironment{longtable}{\fontsize{20pt}{24pt}\selectfont}
\AtBeginEnvironment{tabular}{\fontsize{20pt}{24pt}\selectfont}
\AtBeginEnvironment{equation}{\fontsize{20pt}{24pt}\selectfont}
\AtBeginEnvironment{equation*}{\fontsize{20pt}{24pt}\selectfont}
\AtBeginEnvironment{align}{\fontsize{20pt}{24pt}\selectfont}
\AtBeginEnvironment{align*}{\fontsize{20pt}{24pt}\selectfont}
\AtBeginEnvironment{itemize}{\fontsize{20pt}{24pt}\selectfont}
\AtBeginEnvironment{enumerate}{\fontsize{20pt}{24pt}\selectfont}
\AtBeginEnvironment{description}{\fontsize{20pt}{24pt}\selectfont}
\AtBeginDocument{\usebeamerfont{normal text}}
\newlength{\AFSlideTableWidth}
% The fixed footer reservation is calibrated with the semantic seven-line regular-body budget.
\setbeamertemplate{footline}{%
  \leavevmode\vbox{%
    \hbox{%
      \begin{beamercolorbox}[wd=\paperwidth,ht=16pt,dp=3pt,center]{author in head/foot}%
        \fontsize{16pt}{19pt}\selectfont Untagged PDF derivative \textbar\ \href{../web/index.html}{HTML reader}%
      \end{beamercolorbox}%
    }%
    \vskip6pt%
  }%
}
% Accessible footer reserved height: 25pt.

\newtheorem{proposition}{Proposition}
\newtheorem{hypothesis}{Hypothesis}
\newtheorem{remark}[theorem]{Remark}
\newtheorem{axiom}{Axiom}
\newtheorem{property}{Property}
\usepackage{bookmark}
\IfFileExists{xurl.sty}{\usepackage{xurl}}{} % add URL line breaks if available
\urlstyle{same}
% fallback for those not using the hyperref driver hyperxmp:
\makeatletter
\@ifundefined{xmpquote}{\newcommand{\xmpquote}[1]{#1}}{}
\makeatother
\hypersetup{
  hidelinks,
  pdfcreator={LaTeX via pandoc}}

\author{\texorpdfstring{}{}}
\date{}

\begin{document}

\section{Contributions and evidence boundaries}\label{sec:contributions}

\begin{frame}{Contributions and evidence boundaries}
We make ten contributions, each paired with a theorem, figure, table, or
generated token and with an explicit boundary on what the evidence
establishes.
\end{frame}

\begin{frame}{Contributions and evidence\ldots{} (part 2)}
They fall into four groups: the first two build the core and its
recovery contract; the next two report and statistically qualify the
contaminated-consensus result; the fifth and sixth make the robustness
axes explicit and add the objective-backed server rule; and the
remaining four are scoped extensions ---
\end{frame}

\begin{frame}{Contributions and evidence\ldots{} (part 3)}
parameter recovery, the tempered aggregation family, an aggregation-API
transfer, and a disjoint-observation communication test.
\end{frame}

\begin{frame}{1. A discrete-categorical FedGVI core}
\protect\phantomsection\label{a-discrete-categorical-fedgvi-core}
A typed, deterministic, pure-NumPy/SciPy reimplementation of the FedGVI
(Mildner et al. 2025) generalized-Bayes primitives --- divergences,
bounded robust losses, the generalized posterior, the cavity/factor
algebra, and robust aggregation ---
\end{frame}

\begin{frame}{1. A discrete-categorical\ldots{} (part 2)}
in the discrete-categorical setting that active inference (Da Costa et
al. 2020) uses.
\end{frame}

\begin{frame}{1. A discrete-categorical\ldots{} (part 3)}
The objective (eq.~2) and its closed-form
tempered-softmax solution are stated in Definition
1 and tested by the recovery-limit probes. The
whole core is zero-mock and reproducible (Peng 2011).
\end{frame}

\begin{frame}{2. A recovery-tested connection between categorical
pooling and robust Bayes}
\protect\phantomsection\label{a-recovery-tested-connection-between-categorical-pooling-and-robust-bayes}
\end{frame}

\begin{frame}{2. A recovery-tested\ldots{} (part 2)}
A recovery certificate shows that the client KL/negative-log-likelihood
loss limits recover Bayes and the server's zero-robustness branch
recovers the project log-linear pool.
\end{frame}

\begin{frame}{2. A recovery-tested\ldots{} (part 3)}
Under the explicit shared-support, posterior-log-potential, and
fixed-weight assumptions in sec.~5.8, that pool
specializes Friston et al.'s Eq. 7 message-combination term (Friston et
al. 2024), not the complete source protocol.
\end{frame}

\begin{frame}{2. A recovery-tested\ldots{} (part 4)}
Three recovery limits are pinned to bit-level residuals. The bounded
losses recover NLL/Bayes (Corollary 6 +
Proposition 4; \(\le 0\) and \(\le 0\)
maximum residual, eq.~11).
\end{frame}

\begin{frame}{2. A recovery-tested\ldots{} (part 5)}
The Rényi divergence recovers KL (Lemma 3;
\(\le 0\) residual, eq.~6). The server-side
reweighting pool recovers the naive pool in its trusting limit (Theorem
5; \(\le 0\) residual,
eq.~10).
\end{frame}

\begin{frame}{2. A recovery-tested\ldots{} (part 6)}
Robustness is thereby a tested recovery-limit extension, not a
replacement.
\end{frame}

\begin{frame}{3. End-to-end evaluation of robust federated active
inference}
\protect\phantomsection\label{end-to-end-evaluation-of-robust-federated-active-inference}
Three worked categorical source-mechanism analogues cover communicating
colonies reaching lower free energy
(sec.~7.2), Dirichlet language acquisition
(sec.~7.3),
\end{frame}

\begin{frame}{3. End-to-end evaluation of\ldots{} (part 2)}
and a configured Bayesian-model- reduction pruning sign control (Friston
and Penny 2011) (sec.~7.4).
\end{frame}

\begin{frame}{3. End-to-end evaluation of\ldots{} (part 3)}
A contaminated-sentinel robustness sweep
(sec.~7.5) shows the naive pool degrading while
at least one server-side robust member clears the configured threshold
at the most severe swept rate.
\end{frame}

\begin{frame}{4. A statistically qualified server-side contrast}
\protect\phantomsection\label{a-statistically-qualified-server-side-contrast}
The ``robust beats naive'' conclusion for the declared contamination
rate is produced only by a matched-pairs Wilcoxon signed-rank test
(Wilcoxon 1945).
\end{frame}

\begin{frame}{4. A statistically qualified\ldots{} (part 2)}
Its multiplicity correction uses Benjamini--Hochberg FDR (Benjamini and
Hochberg 1995) across the divergence family.

We report the contrast with bootstrap confidence intervals (Efron and
Tibshirani 1993) and observed-effect design-power planning.
\end{frame}

\begin{frame}{4. A statistically qualified\ldots{} (part 3)}
Across 960 paired trials, the headline display method (RKL; tied set:
RKL, AR, beta, rcce) reaches accuracy 0.9829 against the naive pool's
0.9021 at the verdict rate.
\end{frame}

\begin{frame}{4. A statistically qualified\ldots{} (part 4)}
At that rate, \(q = 1.11 \times 10^{-158}\), and the
rank-biserial-derived \(d\)-equivalent is saturated (r=+1). The
predeclared selection rule is largest positive rank-biserial
effect\_size; stable method order tie-break; the method with the largest
paired mean difference is AR. Every headline number is a generated
token.
\end{frame}

\begin{frame}{5. An explicit accounting of the robustness axes}
\protect\phantomsection\label{an-explicit-accounting-of-the-robustness-axes}
The map in sec.~3.3 separates the client-side
per-agent update, which inherits FedGVI's bounded-influence result under
its stated source assumptions, from both server rules.
\end{frame}

\begin{frame}{5. An explicit accounting of\ldots{} (part 2)}
The sharp server-side reweighting heuristic has a recovery limit as its
positive formal property and a scoped no-go result for its declared
separable objective class.
\end{frame}

\begin{frame}{5. An explicit accounting of\ldots{} (part 3)}
The conservative variational server rule is objective-backed and has a
redescending effective-weight update, but is not the accuracy-maximizer.
Downstream users are told exactly which result is theoretically backed
and which is a labeled heuristic.
\end{frame}

\begin{frame}[fragile]{6. An objective-backed server aggregator with
redescending weights}
\protect\phantomsection\label{an-objective-backed-server-aggregator-with-redescending-weights}
We derive an aggregation free energy eq.~12 whose
exact block updates define the \texttt{variational\_aggregate} rule
(sec.~5.8.2, sec.~12).
\end{frame}

\begin{frame}{6. An objective-backed\ldots{} (part 2)}
Each exact block update monotonically decreases a stated objective, and
a converged fixed point is coordinatewise stationary. The implementation
keeps the lowest observed objective among converged configured starts,
\end{frame}

\begin{frame}{6. An objective-backed\ldots{} (part 3)}
or reports the best unfinished trace as non-converged.
\end{frame}

\begin{frame}{6. An objective-backed\ldots{} (part 4)}
The rule recovers the standard log-linear pool in the trusting limit
(eq.~10) and --- unlike the sharp heuristic ---
carries a proven raw effective-weight bound
(fig.~16, fig.~17).
\end{frame}

\begin{frame}{6. An objective-backed\ldots{} (part 5)}
The honest cost is conservatism: it is a maximum-entropy-biased
consensus and trades peak point-accuracy for that control, so it
complements rather than replaces the sharp heuristic of contribution 5.
\end{frame}

\begin{frame}{7. Executed finite-grid acuity-recovery experiment}
\protect\phantomsection\label{executed-finite-grid-acuity-recovery-experiment}
At each value in the 0.60, 0.70, 0.80, 0.90 acuity grid, the study
generates 200 synthetic observations in each of 960 trials. It selects
acuity by marginal-likelihood grid search over the declared finite grid.
\end{frame}

\begin{frame}{7. Executed finite-grid\ldots{} (part 2)}
The observed mean absolute error is 0.0232 with \(R^2 = 0.9999\)
(fig.~31). Acuity-by-colony-size behavior
belongs to the separate sensitivity study; it is not a
parameter-recovery result.
\end{frame}

\begin{frame}{8. The tempered aggregation family}
\protect\phantomsection\label{the-tempered-aggregation-family}
A one-parameter \(\lambda>0\) generalization of the variational
aggregate (sec.~12.5) is the objective
\end{frame}

\begin{frame}{8. The tempered\ldots{} (part 2)}
\begin{equation}\tag{1}\protect\phantomsection\label{eq:contribution-tempered-family}{
\begin{aligned}
F_\lambda(q,a)
&= \sum_n a_n\,\mathrm{CE}(q,q_n)-\lambda H(q)\\
&\quad +\frac{1}{c}\,\mathrm{KL}_{\mathrm{gen}}(a\Vert w).
\end{aligned}
}\end{equation}
\end{frame}

\begin{frame}[fragile]{8. The tempered\ldots{} (part 3)}
In eq.~1, at \(\lambda = 1.0\), the
temperature is unity and the objective reduces to the standard
variational aggregate bit-for-bit. Lower \(\lambda\) sharpens the
variational \(q\)-block toward a maximizing state; it does not
algebraically recover \texttt{robust\_aggregate}.
\end{frame}

\begin{frame}{8. The tempered\ldots{} (part 4)}
The raw effective-weight update and its bound are preserved for all
\(\lambda\). The full derivation is in sec.~12.5.
\end{frame}

\begin{frame}[fragile]{9. An aggregation-API transfer demonstration}
\protect\phantomsection\label{an-aggregation-api-transfer-demonstration}
The same \texttt{robust\_aggregate} API that governs the POMDP studies
is exercised unchanged with one deterministic MLP trained with the
density-power \(\beta\)-loss (16 hidden units, \(\beta=0.5\);
\end{frame}

\begin{frame}{9. An aggregation-API\ldots{} (part 2)}
generalized variational inference with a point-mass variational family)
as the per-client model.
\end{frame}

\begin{frame}[fragile]{9. An aggregation-API\ldots{} (part 3)}
This supports portability of the server API to one additional model
class (sec.~7.6) when the optional \texttt{torch}
extra is installed (sec.~5.26). Without it, the
MLP run is skipped and its tokens render accordingly.
\end{frame}

\begin{frame}{10. Communication contrast under disjoint observations}
\protect\phantomsection\label{communication-contrast-under-disjoint-observations}
A multi-agent extension (sec.~14.7.1) in
\end{frame}

\begin{frame}{10. Communication contrast\ldots{} (part 2)}
which 3 agents each observe a 2-slot disjoint window shows that belief
sharing materially improves over isolated-agent accuracy in the declared
configuration.
\end{frame}

\begin{frame}{10. Communication contrast\ldots{} (part 3)}
Isolated agents clear the 0.167 chance baseline but stay far below the
communicating consensus, which itself remains well short of full
accuracy.
\end{frame}

\begin{frame}{10. Communication contrast\ldots{} (part 4)}
Across 128 seeds, isolated accuracy is 0.326 versus communicating 0.493,
a reproducible margin under the declared matched-seed comparison
(Wilcoxon \(p = 9.35 \times 10^{-23}\)). This is evidence for the
configured disjoint-observation protocol, not a universal communication
theorem.
\end{frame}

\begin{frame}{10. Communication contrast\ldots{} (part 5)}
The remainder of the paper proceeds as follows. sec.~5
develops the FedGVI core and the recovery limits;
sec.~6 states the numbered recovery theorems and the
expected-free-energy identity;
sec.~5.22 fixes the configuration;
\end{frame}

\begin{frame}{10. Communication contrast\ldots{} (part 6)}
sec.~7 reports the 9 studies, beginning with the
recovery checks (sec.~7.1).
\end{frame}

\begin{frame}{10. Communication contrast\ldots{} (part 7)}
sec.~8 and sec.~9 synthesize; and
sec.~10 and sec.~8.15
document determinism, scope, limitations, and the standing of each
robustness axis.
\end{frame}

\end{document}
